% --- Archon physics formalization source begin ---
% archon:physics
% archon:covers PhyXMiniProblems/problem_phyx_mini_0138.lean
% archon:source-report reports/phyx_mini/problem_phyx_mini_0138.source.json

\chapter{Physics problem phyx\_mini\_0138}
\label{ch:PhyXMiniProblems_problem_phyx_mini_0138}

\paragraph{Problem source.}
\#\# Physical scenario

Light traveling in air strikes a flat piece of uniformly thick glass at an incident angle of \textbackslash{}(60.0\textasciicircum{}\textbackslash{}circ\textbackslash{}), as shown in figure.The index of refraction of the glass is \textbackslash{}(1.50\textbackslash{}).

\#\# Question

What is the angle \textbackslash{}(	heta\_B\textbackslash{}) at which the ray emerges from the glass?

\#\# Answer choices

- A: A: 0.666

- B: B: 0.966

- C: C: 0.866

- D: D: 0.766

\#\# Figure caption (auxiliary; use the image as primary evidence)

The image depicts a diagram illustrating the refraction of light as it passes through a glass slab. 

Key elements include:

1. **Medium Labels**: The diagram is divided into three sections: "Air" on the left and right, and "Glass" in the middle.

2. **Incident Ray**: A solid magenta line labeled "Ray from object" enters the glass at an angle from the left side. The angle of incidence is marked as 60.0°.

3. **Refracted Ray**: Upon entering the glass, the ray bends and is labeled with angle \textbackslash{}(\textbackslash{}theta\_A\textbackslash{}).

4. **Emerging Ray**: The ray exits the glass and refracts again into the air, labeled with angle \textbackslash{}(\textbackslash{}theta\_B\textbackslash{}).

5. **Dashed Line**: A magenta dashed line indicates the position of the "Image" (where the object appears to be) when viewed from above through the glass.

6. **Angles and Labels**: The angles \textbackslash{}(\textbackslash{}theta\_A\textbackslash{}) and \textbackslash{}(\textbackslash{}theta\_B\textbackslash{}) are depicted at the boundaries of the glass.

This diagram illustrates the concept of refraction and how light bends when passing through different media.

\#\# Dataset metadata

Geometrical Optics; Physical Model Grounding Reasoning, Spatial Relation Reasoning

\paragraph{Recorded answer/context.}
C: C: 0.866

\paragraph{Figure/image path.}
/root/proposal\_for\_physic/science-mango/phyx\_mini\_run/phyx\_data/test\_image/138.png

\paragraph{Formalization target.}
create a compiling Lean file with sorry bodies at `PhyXMiniProblems/problem\_phyx\_mini\_0138.lean`. The Lean declarations must preserve the physical quantities, dimensions or dimensional roles, figure labels, governing-law hypotheses, and final relation expressed by this problem.
Use Mathlib/Physlib names found through LeanExplore where available. If a physics API is missing, introduce faithful local abstractions rather than scalar placeholder aliases.

\begin{theorem}[Physics formalization target]
\label{thm:physics:phyx_mini_0138:target}
\lean{PhyXMiniProblems.ProblemPhyXMini0138.problem_phyx_mini_0138}
\uses{def:physics:phyx-mini-0138:phyxminiproblems-problemphyxmini0138-parallelglassslabsetup, def:physics:phyx-mini-0138:phyxminiproblems-problemphyxmini0138-hasphysicalparameters, def:physics:phyx-mini-0138:phyxminiproblems-problemphyxmini0138-matchesproblemreadouts, def:physics:phyx-mini-0138:phyxminiproblems-problemphyxmini0138-matchessuppliedfigure, def:physics:phyx-mini-0138:phyxminiproblems-problemphyxmini0138-satisfiesplaneparallelgeometry, def:physics:phyx-mini-0138:phyxminiproblems-problemphyxmini0138-satisfiessnellslaw, def:physics:phyx-mini-0138:phyxminiproblems-problemphyxmini0138-isuniquematchingdisplayedsine}
The assigned autoformalize agent should translate this physics problem into Lean declarations in the covered file, with theorem and lemma proof bodies written as `by sorry`.
\end{theorem}
\begin{proof}
This is an autoformalization task, not a proof task. Produce statements that can later be proved by the physics prover without weakening the physical model.
\end{proof}

% --- Archon declaration topology begin ---
\section{Declaration topology}
\begin{definition}[Length Quantity]
  \label{def:physics:phyx-mini-0138:phyxminiproblems-problemphyxmini0138-lengthquantity}
  \lean{PhyXMiniProblems.ProblemPhyXMini0138.LengthQuantity}
  A unit-independent physical quantity carrying the dimension of length
\end{definition}
\begin{proof}
  This is the declared model definition of Length Quantity.
\end{proof}
\begin{definition}[length In Meters]
  \label{def:physics:phyx-mini-0138:phyxminiproblems-problemphyxmini0138-lengthinmeters}
  \lean{PhyXMiniProblems.ProblemPhyXMini0138.lengthInMeters}
  \uses{def:physics:phyx-mini-0138:phyxminiproblems-problemphyxmini0138-lengthquantity}
  Meter readout used only to express positivity of the slab thickness
\end{definition}
\begin{proof}
  This is the declared model definition of length In Meters.
\end{proof}
\begin{definition}[Diagram Point]
  \label{def:physics:phyx-mini-0138:phyxminiproblems-problemphyxmini0138-diagrampoint}
  \lean{PhyXMiniProblems.ProblemPhyXMini0138.DiagramPoint}
  A point in the two-dimensional plane of the source ray diagram
\end{definition}
\begin{proof}
  This is the declared model definition of Diagram Point.
\end{proof}
\begin{definition}[Diagram Ray]
  \label{def:physics:phyx-mini-0138:phyxminiproblems-problemphyxmini0138-diagramray}
  \lean{PhyXMiniProblems.ProblemPhyXMini0138.DiagramRay}
  An oriented direction in the diagram, modulo positive rescaling
\end{definition}
\begin{proof}
  This is the declared model definition of Diagram Ray.
\end{proof}
\begin{definition}[Medium Kind]
  \label{def:physics:phyx-mini-0138:phyxminiproblems-problemphyxmini0138-mediumkind}
  \lean{PhyXMiniProblems.ProblemPhyXMini0138.MediumKind}
  Physical kinds of optical medium occurring in the problem
\end{definition}
\begin{proof}
  This is the declared model definition of Medium Kind.
\end{proof}
\begin{definition}[Medium Region]
  \label{def:physics:phyx-mini-0138:phyxminiproblems-problemphyxmini0138-mediumregion}
  \lean{PhyXMiniProblems.ProblemPhyXMini0138.MediumRegion}
  The three regions, ordered from left to right in the source image
\end{definition}
\begin{proof}
  This is the declared model definition of Medium Region.
\end{proof}
\begin{definition}[Slab Face]
  \label{def:physics:phyx-mini-0138:phyxminiproblems-problemphyxmini0138-slabface}
  \lean{PhyXMiniProblems.ProblemPhyXMini0138.SlabFace}
  The entry and exit faces of the glass slab
\end{definition}
\begin{proof}
  This is the declared model definition of Slab Face.
\end{proof}
\begin{definition}[Ray Label]
  \label{def:physics:phyx-mini-0138:phyxminiproblems-problemphyxmini0138-raylabel}
  \lean{PhyXMiniProblems.ProblemPhyXMini0138.RayLabel}
  The solid physical rays and the dashed apparent-image extension
\end{definition}
\begin{proof}
  This is the declared model definition of Ray Label.
\end{proof}
\begin{definition}[Ray Style]
  \label{def:physics:phyx-mini-0138:phyxminiproblems-problemphyxmini0138-raystyle}
  \lean{PhyXMiniProblems.ProblemPhyXMini0138.RayStyle}
  Line styles distinguished in the source image
\end{definition}
\begin{proof}
  This is the declared model definition of Ray Style.
\end{proof}
\begin{definition}[Figure Color]
  \label{def:physics:phyx-mini-0138:phyxminiproblems-problemphyxmini0138-figurecolor}
  \lean{PhyXMiniProblems.ProblemPhyXMini0138.FigureColor}
  The color used for all ray segments in the source image
\end{definition}
\begin{proof}
  This is the declared model definition of Figure Color.
\end{proof}
\begin{definition}[Slab Model]
  \label{def:physics:phyx-mini-0138:phyxminiproblems-problemphyxmini0138-slabmodel}
  \lean{PhyXMiniProblems.ProblemPhyXMini0138.SlabModel}
  The slab idealization stated in the problem
\end{definition}
\begin{proof}
  This is the declared model definition of Slab Model.
\end{proof}
\begin{definition}[Optical Medium]
  \label{def:physics:phyx-mini-0138:phyxminiproblems-problemphyxmini0138-opticalmedium}
  \lean{PhyXMiniProblems.ProblemPhyXMini0138.OpticalMedium}
  \uses{def:physics:phyx-mini-0138:phyxminiproblems-problemphyxmini0138-mediumkind}
  An optical medium together with its dimensionless refractive index. Positivity is kept in HasPhysicalParameters, so this structure stores data rather than silently adding a governing hypothesis
\end{definition}
\begin{proof}
  This is the declared model definition of Optical Medium.
\end{proof}
\begin{definition}[Parallel Glass Slab Setup]
  \label{def:physics:phyx-mini-0138:phyxminiproblems-problemphyxmini0138-parallelglassslabsetup}
  \lean{PhyXMiniProblems.ProblemPhyXMini0138.ParallelGlassSlabSetup}
  \uses{def:physics:phyx-mini-0138:phyxminiproblems-problemphyxmini0138-lengthquantity, def:physics:phyx-mini-0138:phyxminiproblems-problemphyxmini0138-diagrampoint, def:physics:phyx-mini-0138:phyxminiproblems-problemphyxmini0138-diagramray, def:physics:phyx-mini-0138:phyxminiproblems-problemphyxmini0138-mediumregion, def:physics:phyx-mini-0138:phyxminiproblems-problemphyxmini0138-slabface, def:physics:phyx-mini-0138:phyxminiproblems-problemphyxmini0138-raylabel, def:physics:phyx-mini-0138:phyxminiproblems-problemphyxmini0138-raystyle, def:physics:phyx-mini-0138:phyxminiproblems-problemphyxmini0138-figurecolor, def:physics:phyx-mini-0138:phyxminiproblems-problemphyxmini0138-slabmodel, def:physics:phyx-mini-0138:phyxminiproblems-problemphyxmini0138-opticalmedium}
  All physical objects and scalar angle readouts named by the problem and its figure. thetaAAtEntry and thetaAAtExit are kept separate so that equality of the internal angles is supplied by plane-parallel geometry rather than by definition. In particular, thetaB is an unconstrained field here
\end{definition}
\begin{proof}
  This is the declared model definition of Parallel Glass Slab Setup.
\end{proof}
\begin{definition}[Is Physical Refraction Angle]
  \label{def:physics:phyx-mini-0138:phyxminiproblems-problemphyxmini0138-isphysicalrefractionangle}
  \lean{PhyXMiniProblems.ProblemPhyXMini0138.IsPhysicalRefractionAngle}
  A refraction angle represented by its nonnegative acute real lift
\end{definition}
\begin{proof}
  This is the declared model definition of Is Physical Refraction Angle.
\end{proof}
\begin{definition}[Has Physical Parameters]
  \label{def:physics:phyx-mini-0138:phyxminiproblems-problemphyxmini0138-hasphysicalparameters}
  \lean{PhyXMiniProblems.ProblemPhyXMini0138.HasPhysicalParameters}
  \uses{def:physics:phyx-mini-0138:phyxminiproblems-problemphyxmini0138-lengthinmeters, def:physics:phyx-mini-0138:phyxminiproblems-problemphyxmini0138-mediumregion, def:physics:phyx-mini-0138:phyxminiproblems-problemphyxmini0138-parallelglassslabsetup, def:physics:phyx-mini-0138:phyxminiproblems-problemphyxmini0138-isphysicalrefractionangle}
  Positivity and the angular branch appropriate to ordinary air/glass refraction. These conditions select the physically relevant solution of the sine equations without assigning a numerical value to thetaB
\end{definition}
\begin{proof}
  This is the declared model definition of Has Physical Parameters.
\end{proof}
\begin{definition}[Matches Problem Readouts]
  \label{def:physics:phyx-mini-0138:phyxminiproblems-problemphyxmini0138-matchesproblemreadouts}
  \lean{PhyXMiniProblems.ProblemPhyXMini0138.MatchesProblemReadouts}
  \uses{def:physics:phyx-mini-0138:phyxminiproblems-problemphyxmini0138-parallelglassslabsetup}
  Problem-text data: air on both sides, glass in the middle, glass index 1.50, and incident angle 60° = pi/3. The two air regions are modeled as the same medium; no emergent-angle value occurs here
\end{definition}
\begin{proof}
  This is the declared model definition of Matches Problem Readouts.
\end{proof}
\begin{definition}[Matches Supplied Figure]
  \label{def:physics:phyx-mini-0138:phyxminiproblems-problemphyxmini0138-matchessuppliedfigure}
  \lean{PhyXMiniProblems.ProblemPhyXMini0138.MatchesSuppliedFigure}
  \uses{def:physics:phyx-mini-0138:phyxminiproblems-problemphyxmini0138-raylabel, def:physics:phyx-mini-0138:phyxminiproblems-problemphyxmini0138-parallelglassslabsetup}
  Qualitative and geometric information read from the primary image: a flat uniform slab, solid physical ray segments, a dashed backward extension to the apparent image, and magenta coloring. The extension has the emergent ray's direction. None of these readouts prescribes thetaB
\end{definition}
\begin{proof}
  This is the declared model definition of Matches Supplied Figure.
\end{proof}
\begin{definition}[Satisfies Plane Parallel Geometry]
  \label{def:physics:phyx-mini-0138:phyxminiproblems-problemphyxmini0138-satisfiesplaneparallelgeometry}
  \lean{PhyXMiniProblems.ProblemPhyXMini0138.SatisfiesPlaneParallelGeometry}
  \uses{def:physics:phyx-mini-0138:phyxminiproblems-problemphyxmini0138-parallelglassslabsetup}
  Plane-parallel slab geometry. The equal face normals transfer the internal angle at the entry face to the incidence angle at the exit face. This is a geometric law and does not state the requested external emergent angle
\end{definition}
\begin{proof}
  This is the declared model definition of Satisfies Plane Parallel Geometry.
\end{proof}
\begin{definition}[Snell Law At Interface]
  \label{def:physics:phyx-mini-0138:phyxminiproblems-problemphyxmini0138-snelllawatinterface}
  \lean{PhyXMiniProblems.ProblemPhyXMini0138.SnellLawAtInterface}
  \uses{def:physics:phyx-mini-0138:phyxminiproblems-problemphyxmini0138-opticalmedium}
  Snell's law at one interface, n₁ sin(theta₁) = n₂ sin(theta₂). No matching optics declaration was available in Mathlib/PhysLean, so this local predicate records the physical law directly
\end{definition}
\begin{proof}
  This is the declared model definition of Snell Law At Interface.
\end{proof}
\begin{definition}[Satisfies Snells Law]
  \label{def:physics:phyx-mini-0138:phyxminiproblems-problemphyxmini0138-satisfiessnellslaw}
  \lean{PhyXMiniProblems.ProblemPhyXMini0138.SatisfiesSnellsLaw}
  \uses{def:physics:phyx-mini-0138:phyxminiproblems-problemphyxmini0138-parallelglassslabsetup, def:physics:phyx-mini-0138:phyxminiproblems-problemphyxmini0138-snelllawatinterface}
  Snell's law at both air/glass interfaces. The second equation uses the separate internal exit angle; plane-parallel geometry relates it to thetaAAtEntry
\end{definition}
\begin{proof}
  This is the declared model definition of Satisfies Snells Law.
\end{proof}
\begin{definition}[Answer Choice]
  \label{def:physics:phyx-mini-0138:phyxminiproblems-problemphyxmini0138-answerchoice}
  \lean{PhyXMiniProblems.ProblemPhyXMini0138.AnswerChoice}
  Labels of the four dimensionless values printed with the problem
\end{definition}
\begin{proof}
  This is the declared model definition of Answer Choice.
\end{proof}
\begin{definition}[displayed Answer Value]
  \label{def:physics:phyx-mini-0138:phyxminiproblems-problemphyxmini0138-displayedanswervalue}
  \lean{PhyXMiniProblems.ProblemPhyXMini0138.displayedAnswerValue}
  \uses{def:physics:phyx-mini-0138:phyxminiproblems-problemphyxmini0138-answerchoice}
  The dimensionless number printed beside each answer label
\end{definition}
\begin{proof}
  This is the declared model definition of displayed Answer Value.
\end{proof}
\begin{definition}[recorded Dataset Answer]
  \label{def:physics:phyx-mini-0138:phyxminiproblems-problemphyxmini0138-recordeddatasetanswer}
  \lean{PhyXMiniProblems.ProblemPhyXMini0138.recordedDatasetAnswer}
  \uses{def:physics:phyx-mini-0138:phyxminiproblems-problemphyxmini0138-answerchoice}
  The answer label recorded by the source dataset
\end{definition}
\begin{proof}
  This is the declared model definition of recorded Dataset Answer.
\end{proof}
\begin{definition}[Matches Displayed Sine To Nearest Thousandth]
  \label{def:physics:phyx-mini-0138:phyxminiproblems-problemphyxmini0138-matchesdisplayedsinetonearestthousandth}
  \lean{PhyXMiniProblems.ProblemPhyXMini0138.MatchesDisplayedSineToNearestThousandth}
  \uses{def:physics:phyx-mini-0138:phyxminiproblems-problemphyxmini0138-parallelglassslabsetup, def:physics:phyx-mini-0138:phyxminiproblems-problemphyxmini0138-answerchoice, def:physics:phyx-mini-0138:phyxminiproblems-problemphyxmini0138-displayedanswervalue}
  The displayed decimals are dimensionless and choice C is the nearest- thousandth value of sin(thetaB). This explicitly separates the source's decimal choice from the angle itself, whose exact radian value is the primary target
\end{definition}
\begin{proof}
  This is the declared model definition of Matches Displayed Sine To Nearest Thousandth.
\end{proof}
\begin{definition}[Is Unique Matching Displayed Sine]
  \label{def:physics:phyx-mini-0138:phyxminiproblems-problemphyxmini0138-isuniquematchingdisplayedsine}
  \lean{PhyXMiniProblems.ProblemPhyXMini0138.IsUniqueMatchingDisplayedSine}
  \uses{def:physics:phyx-mini-0138:phyxminiproblems-problemphyxmini0138-parallelglassslabsetup, def:physics:phyx-mini-0138:phyxminiproblems-problemphyxmini0138-answerchoice, def:physics:phyx-mini-0138:phyxminiproblems-problemphyxmini0138-matchesdisplayedsinetonearestthousandth}
  A choice is the unique displayed nearest-thousandth sine value
\end{definition}
\begin{proof}
  This is the declared model definition of Is Unique Matching Displayed Sine.
\end{proof}
\begin{lemma}[emergent Angle eq incident Angle]
  \label{lem:physics:phyx-mini-0138:phyxminiproblems-problemphyxmini0138-emergentangle-eq-incidentangle}
  \lean{PhyXMiniProblems.ProblemPhyXMini0138.emergentAngle_eq_incidentAngle}
  \uses{def:physics:phyx-mini-0138:phyxminiproblems-problemphyxmini0138-parallelglassslabsetup, def:physics:phyx-mini-0138:phyxminiproblems-problemphyxmini0138-hasphysicalparameters, def:physics:phyx-mini-0138:phyxminiproblems-problemphyxmini0138-matchesproblemreadouts, def:physics:phyx-mini-0138:phyxminiproblems-problemphyxmini0138-satisfiesplaneparallelgeometry, def:physics:phyx-mini-0138:phyxminiproblems-problemphyxmini0138-satisfiessnellslaw}
  For a plane-parallel slab surrounded by the same medium, the two Snell-law equations cancel the glass refraction and give equal incident and emergent angles on the physical acute branch
\end{lemma}
\begin{proof}
  The conclusion follows from the governing relations and figure readouts named in the statement of emergent Angle eq incident Angle.
\end{proof}
% --- Archon declaration topology end ---
% --- Archon physics formalization source end ---
